\chapter{THẾ HỆ GIỮA: YOLOv5 -- YOLOv8 (2020--2023)}

%% ========== YOLOv5 ==========
\section{YOLOv5 (2020) --- Dân chủ hóa YOLO}
\textbf{Tác giả:} Glenn Jocher / Ultralytics \hfill \textbf{Framework:} PyTorch \hfill \textit{Không có paper}

\subsection{Ý nghĩa lịch sử}
Chuyển từ Darknet (C) sang PyTorch. Release chỉ 2 tháng sau v4, gây tranh cãi vì không có paper. Tuy nhiên trở thành \textbf{YOLO phổ biến nhất} nhờ \texttt{pip install ultralytics}.

\textit{Ví dụ:} Trước v5, muốn dùng YOLO phải compile code C, cấu hình file \texttt{.cfg} phức tạp. v5 đơn giản hóa thành 3 dòng Python:
\begin{verbatim}
    from ultralytics import YOLO
    model = YOLO('yolov5s.pt')
    results = model('anh.jpg')
\end{verbatim}

\subsection{Kiến trúc}
\begin{itemize}
    \item \textbf{Backbone:} CSPDarknet53 (modified), \textbf{C3 module} --- CSP Bottleneck 3 convolutions. C3 chia features thành 2 nhánh: 1 nhánh qua bottleneck, 1 nhánh bypass $\rightarrow$ concat.
    \item \textbf{Neck:} PANet + \textbf{SPPF} (SPP Fast) --- thay 3 MaxPool song song bằng 3 MaxPool \textbf{nối tiếp}. Kết quả tương đương SPP nhưng nhanh hơn $\sim$2$\times$.
    \item \textbf{Head:} Anchor-based, 3 scales. \textbf{AutoAnchor}: dùng k-means + genetic algorithm tự tìm anchors tối ưu cho dataset riêng.
\end{itemize}

\textit{Ví dụ SPPF:} SPP dùng 3 MaxPool kích thước $5\times5$, $9\times9$, $13\times13$ chạy \textbf{cùng lúc} (3 operations). SPPF dùng 3 MaxPool $5\times5$ chạy \textbf{nối tiếp} --- output tương đương nhưng chỉ cần 1 loại kernel, nhanh hơn trên GPU.

\subsection{5 model sizes}
\begin{table}[H]\centering
\caption{YOLOv5 model variants (COCO val2017, input 640)}
\begin{tabular}{lcccl}\toprule
\textbf{Model} & \textbf{Params} & \textbf{FLOPs} & \textbf{mAP50-95} & \textbf{Phù hợp}\\\midrule
v5n & 1.9M & 4.5G & 28.0\% & Mobile, IoT \\
v5s & 7.2M & 16.5G & 37.4\% & Edge devices \\
v5m & 21.2M & 49.0G & 45.4\% & Cân bằng \\
v5l & 46.5M & 109.1G & 49.0\% & Server \\
v5x & 86.7M & 205.7G & 50.7\% & Max accuracy \\\bottomrule
\end{tabular}
\end{table}

\textbf{Đóng góp chính:} SPPF, AutoAnchor, 5 sizes, export 10+ formats, hyperparameter evolution.\\
\textbf{Hạn chế:} Không có paper nên không peer-reviewed; kiến trúc gần giống v4; vẫn anchor-based.

%% ========== YOLOv6 ==========
\section{YOLOv6 (2022) --- Industrial Deployment}
\textbf{Paper:} ``YOLOv6: A Single-Stage Object Detection Framework for Industrial Applications''\\
\textbf{Tác giả:} Meituan Vision AI (dịch vụ giao đồ ăn lớn nhất Trung Quốc)

\subsection{Bối cảnh}
Meituan cần detect vật thể real-time trên hàng triệu đơn hàng mỗi ngày. Yêu cầu: \textbf{cực nhanh} trên GPU server + khả năng quantize INT8 mà không mất accuracy.

\subsection{Re-parameterization --- Kỹ thuật chủ đạo}
Kiến trúc \textbf{khác nhau giữa training và inference}:

\begin{analogybox}
\textbf{Ôn thi:} Khi ôn (training), bạn đọc 5 cuốn sách, ghi chú, làm bài tập (multi-branch). Khi vào phòng thi (inference), bạn chỉ cần 1 tờ tóm tắt (single branch) chứa tinh hoa từ 5 cuốn.
\end{analogybox}

\begin{itemize}
    \item \textbf{Training:} Multi-branch --- Conv $3\times3$ + Conv $1\times1$ + Identity song song $\rightarrow$ học phong phú
    \item \textbf{Inference:} Gộp (re-parameterize) thành 1 Conv $3\times3$ duy nhất $\rightarrow$ nhanh
    \item \textbf{Toán học:} $W_{merged} = W_{3\times3} + \text{pad}(W_{1\times1}) + I$ --- gộp weights bằng phép cộng ma trận
\end{itemize}

\subsection{Kiến trúc chi tiết}
\begin{itemize}
    \item \textbf{Backbone --- EfficientRep:} Thay CSPDarknet bằng RepVGG blocks. Training: multi-branch. Inference: single-path Conv.
    \item \textbf{Neck --- Rep-PAN:} PAN cũng dùng RepVGG blocks $\rightarrow$ inference nhanh hơn.
    \item \textbf{Head --- Efficient Decoupled Head:} Tách cls/reg (giống YOLOX) nhưng giảm 1 Conv layer $\rightarrow$ tiết kiệm computation.
    \item \textbf{Label Assignment:} SimOTA (giai đoạn 1) $\rightarrow$ \textbf{TAL} (giai đoạn 2): Task-Aligned Learning căn chỉnh cls score và localization quality.
    \item \textbf{Loss:} VFL (Varifocal Loss) + SIoU/GIoU + DFL (model lớn).
    \item \textbf{Anchor-free:} Dự đoán trực tiếp 4 distances từ tâm (FCOS-style).
\end{itemize}

\subsection{Quantization Pipeline --- Điểm mạnh nhất}
\begin{enumerate}
    \item \textbf{PTQ} (Post-Training Quantization): Quantize model đã train $\rightarrow$ nhanh nhưng mất accuracy.
    \item \textbf{QAT} (Quantization-Aware Training): Train lại với simulated quantization $\rightarrow$ giữ accuracy.
    \item \textbf{RepOptimizer:} Tối ưu re-param blocks cho quantization.
\end{enumerate}
Kết quả: INT8 chỉ mất \textbf{0.2\% mAP} nhưng tăng \textbf{75\% FPS}.

\subsection{Kết quả}
\begin{table}[H]\centering
\caption{YOLOv6 so sánh (COCO val, T4 GPU TensorRT FP16)}
\begin{tabular}{lcccc}\toprule
\textbf{Model} & \textbf{mAP50-95} & \textbf{FPS} & \textbf{Params} & \textbf{So sánh} \\\midrule
YOLOv5-N & 28.0\% & 602 & 1.9M & --- \\
\textbf{YOLOv6-N} & \textbf{37.5\%} & \textbf{1187} & \textbf{4.7M} & +9.5\%, 2$\times$ nhanh hơn \\
YOLOv5-S & 37.4\% & 349 & 7.2M & --- \\
\textbf{YOLOv6-S} & \textbf{45.0\%} & \textbf{495} & \textbf{18.5M} & +7.6\% \\
YOLOv5-L & 49.0\% & 113 & 46.5M & --- \\
\textbf{YOLOv6-L} & \textbf{52.8\%} & \textbf{116} & \textbf{59.6M} & +3.8\% \\\bottomrule
\end{tabular}
\end{table}

\textbf{Nhận xét:} v6 vượt v5 ở \textbf{mọi scale}, đặc biệt mạnh ở Nano (1187 FPS!). Tuy nhiên ecosystem nhỏ hơn v5 nhiều.

%% ========== YOLOv7 ==========
\section{YOLOv7 (2022) --- SOTA Accuracy}
\textbf{Paper:} ``YOLOv7: Trainable Bag-of-Freebies Sets New State-of-the-Art for Real-Time Object Detectors'' (CVPR 2023)\\
\textbf{Tác giả:} Chien-Yao Wang, Alexey Bochkovskiy, Hong-Yuan Mark Liao

\subsection{E-ELAN (Extended Efficient Layer Aggregation Network)}
Mở rộng ELAN bằng cách tăng \textbf{cardinality} (số nhánh song song) mà \textbf{không phá gradient path gốc}.

\begin{analogybox}
\textbf{Dây chuyền sản xuất:} ELAN giống 1 dây chuyền với 3 công nhân nối tiếp. E-ELAN thêm ``ca kíp'' mới (expand cardinality) nhưng giữ nguyên dây chuyền cũ, rồi \textbf{shuffle + merge} kết quả $\rightarrow$ output phong phú hơn.
\end{analogybox}

\begin{itemize}
    \item \textbf{Expand:} Tăng số computational blocks (4 $\rightarrow$ 8 branches)
    \item \textbf{Shuffle:} Xáo trộn channels giữa các branches
    \item \textbf{Merge:} Gộp features từ tất cả branches
    \item Kết quả: +0.7\% mAP so với ELAN chuẩn
\end{itemize}

\subsection{Planned Re-parameterized Conv}
Phát hiện quan trọng: RepConv (từ RepVGG) \textbf{không nên dùng} khi có identity connection (residual). Kết hợp identity + conv $1\times1$ + conv $3\times3$ gây \textbf{xung đột gradient}.

\textbf{Giải pháp --- RepConvN:} Bỏ identity branch, chỉ giữ conv $1\times1$ + conv $3\times3$. Tăng 0.5\% AP với cùng speed.

\subsection{Coarse-to-Fine Label Assignment}
\begin{itemize}
    \item \textbf{Lead Head} (fine-grained): Head chính, gán nhãn chặt chẽ
    \item \textbf{Auxiliary Head} (coarse-grained): Head phụ, gán nhãn lỏng hơn $\rightarrow$ nhiều positive samples hơn $\rightarrow$ train tốt hơn
    \item Auxiliary head \textbf{chỉ dùng khi training}, bỏ khi inference $\rightarrow$ zero inference cost
\end{itemize}

\begin{analogybox}
\textbf{Giáo viên và trợ giảng:} Lead head giống giáo viên chấm thi nghiêm khắc (chỉ cho điểm khi đúng hoàn toàn). Auxiliary head giống trợ giảng nhẹ nhàng (cho điểm khi gần đúng) $\rightarrow$ sinh viên nhận nhiều feedback $\rightarrow$ học tốt hơn.
\end{analogybox}

\subsection{Kết quả chi tiết}
\begin{table}[H]\centering
\caption{YOLOv7 Base models (V100 GPU)}
\begin{tabular}{lcccc}\toprule
\textbf{Model} & \textbf{Params} & \textbf{FPS} & \textbf{mAP50-95} & \textbf{So sánh}\\\midrule
YOLOv5-L & 46.5M & 99 & 48.8\% & --- \\
YOLOX-L & 54.2M & 68 & 50.0\% & --- \\
PPYOLOE-L & 52.2M & 78 & 51.4\% & --- \\
\textbf{YOLOv7} & \textbf{36.9M} & \textbf{161} & \textbf{51.4\%} & Ít params, nhanh nhất \\
\textbf{YOLOv7-X} & \textbf{71.3M} & \textbf{114} & \textbf{53.1\%} & Vượt PPYOLOE \\\bottomrule
\end{tabular}
\end{table}

\begin{table}[H]\centering
\caption{YOLOv7 Large models --- SOTA}
\begin{tabular}{lccc}\toprule
\textbf{Model} & \textbf{FPS} & \textbf{mAP50-95}\\\midrule
YOLOv7-W6 & 84 & 54.9\% \\
YOLOv7-E6 & 56 & 55.9\% \\
\textbf{YOLOv7-E6E} & \textbf{36} & \textbf{56.8\%} \\
SWIN-L Cascade R-CNN & 12 & 53.9\% \\
ConvNeXt-XL Cascade R-CNN & --- & 55.2\% \\\bottomrule
\end{tabular}
\end{table}

\textbf{Kết luận:} YOLOv7-E6E (56.8\%) vượt tất cả --- cả Transformer-based (SWIN-L). SOTA accuracy cho đến khi YOLO26 ra đời.

%% ========== YOLOv8 ==========
\section{YOLOv8 (2023) --- Đa năng nhất}
\textbf{Tác giả:} Ultralytics \hfill \textit{Không có paper}

\subsection{3 thay đổi kiến trúc lớn so với v5}
\begin{enumerate}
    \item \textbf{Anchor-Free:} Bỏ hoàn toàn anchor boxes. Dự đoán \textbf{center point + 4 distances} (khoảng cách từ tâm đến 4 cạnh box).
    
    \textit{Ví dụ:} Anchor-based giống chọn khung ảnh có sẵn (S/M/L) rồi chỉnh. Anchor-free giống vẽ tự do --- đứng tại tâm vật, đo 4 hướng lên/xuống/trái/phải.
    
    \item \textbf{Decoupled Head:} Tách classification và regression thành 2 branches hoàn toàn riêng biệt.
    
    \textit{Ví dụ:} Thay vì 1 bác sĩ vừa chẩn đoán bệnh vừa chỉ vị trí đau, chia thành 2 bác sĩ chuyên môn --- mỗi người làm tốt hơn.
    
    \item \textbf{DFL (Distribution Focal Loss):} Thay vì dự đoán 1 giá trị cho mỗi tọa độ box, dự đoán \textbf{phân phối xác suất} trên 16 bins.
    
    \textit{Ví dụ:} Thay vì nói ``cạnh trái cách tâm 25px'' (1 số), DFL nói ``xác suất 10\% ở 24px, 70\% ở 25px, 20\% ở 26px'' $\rightarrow$ phản ánh \textbf{mức độ không chắc chắn}, localization chính xác hơn.
\end{enumerate}

\subsection{C2f Module}
Kế thừa C3 (v5) + cảm hứng từ ELAN (v7):
\begin{itemize}
    \item C3: chỉ dùng output cuối cùng của bottleneck chain
    \item \textbf{C2f: concat output của TẤT CẢ bottlenecks} $\rightarrow$ gradient flow phong phú hơn, nhiều scale features hơn
\end{itemize}

\subsection{Multi-task --- 6 tác vụ trong 1 framework}
\begin{table}[H]\centering
\caption{6 tác vụ YOLOv8 hỗ trợ}
\begin{tabular}{lll}\toprule
\textbf{Tác vụ} & \textbf{Output} & \textbf{Ví dụ ứng dụng} \\\midrule
Detection & Bounding boxes & Phát hiện xe, người \\
Segmentation & Pixel-level masks & Tách nền ảnh \\
Classification & Class label & Phân loại sản phẩm \\
Pose Estimation & Keypoints (17 điểm) & Phân tích tư thế \\
OBB & Rotated boxes & Ảnh vệ tinh, tàu thuyền \\
Tracking & Box + Track ID & Đếm người, theo dõi xe \\\bottomrule
\end{tabular}
\end{table}

\subsection{Kết quả}
\begin{table}[H]\centering
\caption{YOLOv8 so sánh (COCO val2017, input 640)}
\begin{tabular}{lcccc}\toprule
\textbf{Model} & \textbf{Params} & \textbf{FLOPs} & \textbf{mAP50-95} & \textbf{So với v5}\\\midrule
v8n & 3.2M & 8.7G & 37.3\% & v5n: 28.0\% (+9.3\%) \\
v8s & 11.2M & 28.6G & 44.9\% & v5s: 37.4\% (+7.5\%) \\
v8m & 25.9M & 78.9G & 50.2\% & v5m: 45.4\% (+4.8\%) \\
v8l & 43.7M & 165.2G & 52.9\% & v5l: 49.0\% (+3.9\%) \\
v8x & 68.2M & 257.8G & 53.9\% & v5x: 50.7\% (+3.2\%) \\\bottomrule
\end{tabular}
\end{table}

\textbf{Lưu ý quan trọng:} v8-X (53.9\%) \textbf{thua} v7-E6E (56.8\%) về mAP thuần túy! Tuy nhiên v8 thắng ở: ecosystem ($\star\star\star\star\star$), multi-task (6 tác vụ), API thống nhất.

\textbf{Nếu cần accuracy tối đa $\rightarrow$ v7. Nếu cần đa năng + dễ dùng $\rightarrow$ v8.}
